\documentclass{article}
\usepackage{amsmath,amssymb,amsthm,enumitem}
\usepackage{algorithm,algorithmic}
\usepackage{hyperref}
\newtheorem{theorem}{Theorem}
\newtheorem{lemma}{Lemma}
\newtheorem{assumption}{Assumption}
\newcommand{\E}{\mathbb{E}}
\newcommand{\argmin}{\operatorname*{arg\,min}}

\begin{document}

\section*{Frank--Wolfe Algorithm (Conditional Gradient)}

\section*{Mathematical Formulation}
Let $f:\mathbb{R}^d\to\mathbb{R}$ be a convex and $L$-smooth function and let $\mathcal{C}\subset\mathbb{R}^d$ be a non‑empty compact convex set with diameter 
\[
D \;:=\; \max_{x,y\in\mathcal{C}}\|x-y\|_2 .
\]
The Frank–Wolfe method generates a sequence $\{x_k\}_{k\ge 0}$ as follows:

\begin{enumerate}[label=\alph*)]
\item \textbf{Linear minimization oracle (LMO)}  
\[
s_k \;=\; \argmin_{s\in\mathcal{C}} \langle \nabla f(x_k),\, s\rangle .
\tag{1}
\]
\item \textbf{Step‑size rule} (standard choice)  
\[
\gamma_k \;=\; \frac{2}{k+2}, \qquad k=0,1,2,\dots
\tag{2}
\]
\item \textbf{Convex combination update}  
\[
x_{k+1} \;=\; (1-\gamma_k)\,x_k \;+\; \gamma_k\, s_k .
\tag{3}
\]
\end{enumerate}
The algorithm stops when a prescribed maximum number of iterations $T$ is reached or when a duality gap
\[
g(x_k) \;:=\; \langle \nabla f(x_k),\, x_k - s_k\rangle
\]
falls below a tolerance.

\section*{Pseudocode}
\begin{algorithm}[H]
\caption{Frank–Wolfe (Conditional Gradient)}\label{alg:FW}
\begin{algorithmic}[1]
\REQUIRE Convex $L$‑smooth $f$, compact convex $\mathcal{C}$, initial point $x_0\in\mathcal{C}$, max iter $T$.
\FOR{$k=0$ \TO $T-1$}
    \STATE Compute gradient $g_k \gets \nabla f(x_k)$.
    \STATE \textbf{LMO:} $s_k \gets \argmin_{s\in\mathcal{C}} \langle g_k, s\rangle$.
    \STATE Choose step size $\gamma_k \gets \dfrac{2}{k+2}$.
    \STATE Update $x_{k+1} \gets (1-\gamma_k)x_k + \gamma_k s_k$.
    \STATE (Optional) compute duality gap $g(x_k)=\langle g_k, x_k-s_k\rangle$ and stop if $g(x_k)\le\varepsilon$.
\ENDFOR
\RETURN $x_T$.
\end{algorithmic}
\end{algorithm}

\section*{Dry Run Example}
Consider the univariate problem
\[
f(w) = (w-3)^2 ,\qquad \mathcal{C}=[0,5],\qquad w_0 = 0,
\]
and fix a constant step size $\gamma=0.1$ (instead of the standard $\frac{2}{k+2}$) for illustration.  
The gradient is $\nabla f(w)=2(w-3)$.

\begin{center}
\begin{tabular}{c|c|c|c|c}
\textbf{Iter} & $w_k$ & $\nabla f(w_k)$ & $s_k$ (LMO) & $w_{k+1}= (1-\gamma)w_k+\gamma s_k$\\ \hline
0 & $0$ & $-6$ & $\argmin_{s\in[0,5]}(-6)s = 5$ & $0.9\cdot0+0.1\cdot5 = 0.5$\\
1 & $0.5$ & $-5$ & $\argmin_{s\in[0,5]}(-5)s = 5$ & $0.9\cdot0.5+0.1\cdot5 = 0.95$\\
2 & $0.95$ & $-4.1$ & $\argmin_{s\in[0,5]}(-4.1)s = 5$ & $0.9\cdot0.95+0.1\cdot5 = 1.355$\\
\end{tabular}
\end{center}
Corresponding objective values:
\[
f(0)=9,\quad f(0.5)=6.25,\quad f(0.95)=4.2025,\quad f(1.355)\approx 2.705.
\]
The iterates move towards the minimiser $w^\star=3$, illustrating the Frank–Wolfe dynamics.

\newpage
\section*{Assumptions}
\begin{assumption}[Convexity and Smoothness]\label{ass:conv-smooth}
$f$ is convex and differentiable on an open set containing $\mathcal{C}$, and there exists $L>0$ such that
\[
\|\nabla f(x)-\nabla f(y)\|_2 \le L\|x-y\|_2,\qquad \forall x,y\in\mathcal{C}.
\tag{A1}
\]
\end{assumption}
\begin{assumption}[Compact Feasible Set]\label{ass:compact}
$\mathcal{C}$ is compact and convex with finite diameter $D$ defined above.
\tag{A2}
\end{assumption}
\begin{assumption}[Exact Linear Minimization Oracle]\label{ass:lmo}
At each iteration the LMO returns an exact solution $s_k$ of \eqref{1}.
\tag{A3}
\end{assumption}

\section*{Main Convergence Theorem}
\begin{theorem}[Primal Gap $O(1/k)$]\label{thm:rate}
Let $\{x_k\}$ be generated by the Frank–Wolfe algorithm with step size $\gamma_k=2/(k+2)$. Under Assumptions~\ref{ass:conv-smooth}–\ref{ass:lmo} the primal suboptimality satisfies
\[
f(x_k)-f(x^\star) \;\le\; \frac{2L D^{2}}{k+2},\qquad \forall k\ge 0,
\tag{4}
\]
where $x^\star\in\arg\min_{x\in\mathcal{C}}f(x)$.
\end{theorem}
A direct corollary is that the duality gap $g(x_k)$ obeys the same $O(1/k)$ bound.

\section*{Theoretical Properties}
\begin{itemize}
\item \textbf{Stability.} The iterates remain in $\mathcal{C}$ for all $k$ because $x_{k+1}$ is a convex combination of two points of $\mathcal{C}$.
\item \textbf{Hyper‑parameter sensitivity.} The standard schedule $\gamma_k=2/(k+2)$ guarantees the $O(1/k)$ rate without tuning. Constant step sizes lead to a linear convergence neighbourhood that depends on $\gamma$ and the curvature constant.
\item \textbf{Comparison to projected gradient.} Frank–Wolfe avoids costly Euclidean projections; the price is a slower $O(1/k)$ rate compared with $O(1/k)$ for projected gradient under the same smoothness but with potentially larger constants.
\item \textbf{Curvature constant.} Defining
\[
C_f \;:=\; \sup_{\substack{x,s\in\mathcal{C}\\ \gamma\in[0,1]}} \frac{2}{\gamma^{2}}
\Bigl(f\bigl(x+\gamma(s-x)\bigr)-f(x)-\gamma\langle\nabla f(x),s-x\rangle\Bigr),
\tag{5}
\]
one can replace $2LD^{2}$ in \eqref{4} by $C_f$, which may be substantially smaller.
\end{itemize}

\section*{Discussion}
The Frank–Wolfe method is particularly attractive when $\mathcal{C}$ admits a cheap LMO (e.g., simplex, trace‑norm ball, matroid polytopes). The $O(1/k)$ rate is optimal for the class of smooth convex functions without additional strong convexity or curvature assumptions. Accelerated variants (e.g., AFW, line‑search) can improve constants or achieve $O(1/k^{2})$ under stronger curvature conditions.

\newpage
\section*{Preliminaries (Definitions and Lemmas)}

\begin{lemma}[Descent Lemma]\label{lem:descent}
If $f$ is $L$‑smooth then for any $x,y\in\mathcal{C}$,
\[
f(y) \le f(x) + \langle\nabla f(x),y-x\rangle + \frac{L}{2}\|y-x\|_2^{2}.
\tag{6}
\]
\end{lemma}
\begin{proof}
Standard consequence of $L$‑smoothness via integration of the gradient; see e.g. Nesterov (2004). \qedhere
\end{proof}

\begin{lemma}[Curvature Upper Bound]\label{lem:curvature}
The curvature constant $C_f$ defined in \eqref{5} satisfies $C_f \le L D^{2}$.
\tag{7}
\end{lemma}
\begin{proof}
For any $x,s\in\mathcal{C}$ and $\gamma\in[0,1]$, apply \eqref{6} with $y=x+\gamma(s-x)$:
\[
f(x+\gamma(s-x)) \le f(x) + \gamma\langle\nabla f(x),s-x\rangle + \frac{L}{2}\gamma^{2}\|s-x\|^{2}.
\]
Rearranging yields the term inside the supremum in \eqref{5} is at most $L\|s-x\|^{2}\le L D^{2}$. Multiplying by $2/\gamma^{2}$ gives \eqref{7}. \qedhere
\end{proof}

\section*{Main Proof (Step‑by‑Step Derivation)}
We prove Theorem~\ref{thm:rate}. Throughout the proof we denote $h_k:=f(x_k)-f(x^\star)\ge0$.

\begin{enumerate}
\item[Step 1.] \textbf{Use curvature definition.} By definition of $C_f$ and the update $x_{k+1}=x_k+\gamma_k(s_k-x_k)$,
\[
f(x_{k+1}) \;\le\; f(x_k) + \gamma_k\langle\nabla f(x_k),s_k-x_k\rangle + \frac{\gamma_k^{2}}{2}C_f .
\tag{8}
\]
\item[Step 2.] \textbf{Relate linear oracle to duality gap.} Since $s_k$ minimizes $\langle\nabla f(x_k),s\rangle$ over $\mathcal{C}$,
\[
\langle\nabla f(x_k),s_k-x_k\rangle \;=\; -\,g(x_k),
\tag{9}
\]
where $g(x_k):=\langle\nabla f(x_k),x_k-s_k\rangle\ge0$ is the Frank–Wolfe (duality) gap.
\item[Step 3.] \textbf{Upper bound the gap by primal suboptimality.} Convexity of $f$ gives
\[
f(x^\star) \;\ge\; f(x_k) + \langle\nabla f(x_k),x^\star-x_k\rangle .
\tag{10}
\]
Since $s_k$ is the minimizer of the linear form,
\[
\langle\nabla f(x_k),s_k-x_k\rangle \;\le\; \langle\nabla f(x_k),x^\star-x_k\rangle .
\tag{11}
\]
Combining \eqref{9}–\eqref{11} yields
\[
g(x_k) \;\ge\; f(x_k)-f(x^\star)=h_k .
\tag{12}
\]
\item[Step 4.] \textbf{Insert (12) into (8).} Using (9) and (12),
\[
f(x_{k+1}) \;\le\; f(x_k) - \gamma_k g(x_k) + \frac{\gamma_k^{2}}{2}C_f
\;\le\; f(x_k) - \gamma_k h_k + \frac{\gamma_k^{2}}{2}C_f .
\tag{13}
\]
Rearranging,
\[
h_{k+1} \;\le\; (1-\gamma_k)h_k + \frac{\gamma_k^{2}}{2}C_f .
\tag{14}
\]
\item[Step 5.] \textbf{Plug the standard step size $\gamma_k=\frac{2}{k+2}$.} Then $1-\gamma_k=\frac{k}{k+2}$ and $\gamma_k^{2}=\frac{4}{(k+2)^{2}}$. Substituting into \eqref{14}:
\[
h_{k+1} \;\le\; \frac{k}{k+2}h_k + \frac{2}{(k+2)^{2}}C_f .
\tag{15}
\]
\item[Step 6.] \textbf{Induction to obtain closed form.} We claim
\[
h_k \;\le\; \frac{2C_f}{k+2},\qquad \forall k\ge0 .
\tag{16}
\]
\emph{Base case $k=0$:} $h_0 = f(x_0)-f(x^\star)\le C_f$ (since $C_f\ge0$) and $\frac{2C_f}{2}=C_f$, so (16) holds.

\emph{Inductive step:} Assume (16) holds for some $k$. Using \eqref{15},
\[
h_{k+1} \;\le\; \frac{k}{k+2}\cdot\frac{2C_f}{k+2}+ \frac{2C_f}{(k+2)^{2}}
= \frac{2C_f}{(k+2)^{2}}\bigl(k+1\bigr)
= \frac{2C_f}{(k+3)}\cdot\frac{k+1}{k+2}.
\]
Since $\frac{k+1}{k+2}\le 1$, we obtain
\[
h_{k+1} \le \frac{2C_f}{k+3},
\]
which is precisely (16) with $k\leftarrow k+1$. Thus the claim holds for all $k$.
\item[Step 7.] \textbf{Replace $C_f$ by $L D^{2}$.} By Lemma~\ref{lem:curvature}, $C_f\le L D^{2}$. Inserting into (16) yields
\[
h_k \;\le\; \frac{2L D^{2}}{k+2},
\]
which is exactly the bound \eqref{4}. \qedhere
\end{enumerate}

\section*{Rate Derivation (With Constants)}
The derivation above shows that the primal gap decays as
\[
f(x_k)-f(x^\star) \le \frac{2L D^{2}}{k+2},
\]
where:
\begin{itemize}
\item $L$ is the Lipschitz constant of $\nabla f$;
\item $D$ is the Euclidean diameter of $\mathcal{C}$;
\item the factor $2$ stems from the choice $\gamma_k=2/(k+2)$.
\end{itemize}
If a tighter curvature constant $C_f$ is known, replace $2L D^{2}$ by $C_f$ for a sharper bound.

\section*{Special Cases}
\begin{enumerate}
\item \textbf{Strongly convex $f$.} If $f$ is $\mu$‑strongly convex, the duality gap satisfies $g(x_k)\ge \mu \|x_k-x^\star\|^{2}$, and the $O(1/k)$ rate can be upgraded to linear convergence when a \emph{away‑step} or \emph{pairwise} variant is used.
\item \textbf{Exact line search.} Choosing $\gamma_k = \arg\min_{\gamma\in[0,1]} f\bigl(x_k+\gamma(s_k-x_k)\bigr)$ yields the same $O(1/k)$ rate but with a potentially smaller constant because the optimal $\gamma_k$ satisfies
\[
\gamma_k = \frac{g(x_k)}{C_f}\le \frac{2}{k+2}.
\]
\item \textbf{Polyhedral feasible set.} When $\mathcal{C}$ is a polytope, the Frank–Wolfe iterates are sparse (they lie on a face with at most $k+1$ vertices), which is advantageous for high‑dimensional problems.
\end{enumerate}

\end{document}